


\chapter{Introduction}

\section{Introduction}
Real-time processing of large networks and determining their community structures has a wide variety of applications such as finding clusters in social media networks, optimising resource allocation in real time data processing tasks and more. The typical number of new network edges produced on social media platforms now count in the millions while the size of the graphs themselves can be much higher. This explosion in size of networks over the recent years mean that traditional non-incremental graph clustering algorithms cannot be run when new nodes or edges are added and still yield timely information, which is crucial in extremely dynamic social networks. While there exists a large body of research concerning graph partitioning algorithms, unfortunately there has been little research as to how these algorithms could be adopted or used in an online setting. Additionally, the handful of existing literature concerning real time graph partitioning mostly limits itself to using simple, inexpensive heuristics. Incremental computation through techniques such as Differential Dataflow \cite{mcsherry2013differential} could allow for more advanced and computationally demanding algorithms such as the Louvain \cite{blondel2008fast} modularity optimisation algorithm to be used in online settings.\\
\textbf{Contributions:} The main contributions of this paper are:
\begin{itemize}
    \item Introduction of the viability of using the differential dataflow computational model for the purposes of graph processing. 
    \item Demonstrate how incremental computation could be used to significantly speedup traditionally non-incremental algorithms that perform lots of unnecessary work such as the Louvain modularity optimisation algorithm. 
    \item Demonstrate and explore the limitations of parallelising Louvain and provide potential mitigation strategies where appropriate. 
\end{itemize}


\section{Problem Statement and notation}
Let $G(V,E)$ be a graph such that $V$ represents the set of vertices and $E$ represents the set of edges from one vertex to another. It is important to note that this graph is also undirected, while algorithms have been researched that apply modularity optimisation to directed graphs, we only concern ourselves with undirected graphs for the purposes of this paper. In this graph we allow for edges that connect a vertex to itself and we define this as a ``self-loop'', however, multi-edges are prohibited and in the case where the dataset contains multi-edges, these must be first agglomerated into a single edge by simply adding the costs in each multi-edge. 
A \textit{community} within a graph represents a set of vertices (possibly empty) such that the edges of these vertices display strong density. In practice however, we only concern ourselves with partitioning the vertex set $V$ into some number of \textit{disjoint} non-empty communities. We will use $C(i)$ as notation to indicate the community of a given vertex $i$. 
We will define the adjacency matrix that defines our graph through $A_{ij}$, that is $A_{ij}$ represents the weight from vertex $i$ to $j$. Let $k_i$ denote the weighted degree of a given vertex $k_i$, that is $k_i = \sum_{j}A_{ij}$ and finally we define $m$ to be the summation of all the edges present in the graph of interest, that is $m = \frac{1}{2}\sum_{ij}A_{ij}$, note we half the weight as a result of the doubling of weights present when an undirected graph is represented as an adjacency matrix. Finally we can now define modularity as $Q=\frac{1}{2 m} \sum_{i, j}\left[A_{i j}-\frac{k_{i} k_{j}}{2 m}\right] \delta\left(C(i), C(j)\right)$. Note that the $\delta\left(u,v\right)$ is simply the Kronecker delta function.  

\section{Motivation}
\subsection{Differential Dataflow}
Differential Dataflow \cite{mcsherry2013differential} is an incremental computational model that has shown promise in a variety of graph processing tasks. In their paper, they demonstrated that their computational model was able to solve graph algorithms such as PageRank, Strongly Connected Components with a significant increase in performance over a selection of baselines, the speedups ranged from 1.9x to 1899.6x. Given most greedy modularity optimisation methods require multiple iterations over non-changing, already calculated vertices, it therefore makes sense to warrant an investigation on how Differential Dataflow could be used to reduce the cost of these iterations, additionally when we consider a dynamically changing graph we note that the increase in possible community swaps do not immediately increase with respect to the number of vertices but rather with respect to the degree of the inserted edge, although a swap as a result of this insertion may result in a domino effect of changes later on still. 

\subsubsection{Timely Dataflow}
The open source implementation of Differential Dataflow was also a critical advantage when developing with Differential Dataflow, it provided easy access for distributed processing of data or in our case graphs. This ease of access for distributed computations is provided by the underlying dataflow computation engine called `Timely Dataflow' first introduced in Naiad \cite{murray2013naiad} which states that they provide a framework for low-latency cyclic dataflow computations. This engine ensured that our implementation could run on a cluster, a process or multiple processes with relatively little configuration.

\subsection{Louvain}
One criteria for choosing Louvain modularity optimisation as a candidate for our exploration of incremental, data-parallel graph clustering was simply as a result of its usefulness and widespread use in both academia and in industry. Louvain despite its pitfalls has remained the most commonly used graph clustering algorithm with various open source implementations in distributed dataflow processing frameworks such as Apache Spark \footnote{See \url{https://spark.apache.org/}}. The other criteria for choosing Louvain is simply due to the unnecessary work it performs when iterating over vertices, this provides a place of improvement in terms of performance, additionally as mentioned previously the advantages of incrementally performing Louvain may result in showing that the use of Louvain like methods used for agglomerative clustering could also be used in conjunction with Differential Dataflow. 